% Lösungen Einheit 2 von 2 – Kumulieren gegen eine Schranke: hypergeometrische Einzelwahrscheinlichkeiten aufsummieren und die größte (oder kleinste) Trefferzahl gegen eine Schranke bestimmen (zusammenbau v0.1)
\einheitenkopf{Einheit 2 von 2 · Kumulieren gegen eine Schranke: hypergeometrische Einzelwahrscheinlichkeiten aufsummieren und die größte (oder kleinste) Trefferzahl gegen eine Schranke bestimmen}

\erg{9}{a) $\frac{(4 \text{ über } 2) \cdot (6 \text{ über } 1)}{(10 \text{ über } 3)}$ – günstige durch mögliche Auswahlen, im Zähler Treffer und Nieten. \quad b) $P(X = 0) \approx 0{,}1667$; $P(X = 1) \approx 0{,}5000$; $P(X = 2) \approx 0{,}3000$; $P(X \le 1) \approx 0{,}6667 < 0{,}7$, $P(X \le 2) \approx 0{,}9667 \ge 0{,}7$; größte Trefferzahl $k = 1$ \quad c) $P(X = 0) \approx 0{,}0372$, $P(X = 1) \approx 0{,}2008$, $P(X = 2) \approx 0{,}3651$; $P(X \le 1) \approx 0{,}2380 < 0{,}4$, $P(X \le 2) \approx 0{,}6031 \ge 0{,}4$; also $n = 1$. Dass $P(X = 2) < 0{,}4$ ist, zählt nicht – verglichen wird die Summe.}
\erg{10}{a) Fehler: Nina kumuliert binomial, obwohl ohne Zurücklegen entnommen wird. Richtig: $P(X \le 0) = \frac{1}{6} \approx 0{,}1667 < 0{,}66$, $P(X \le 1) = \frac{2}{3} \approx 0{,}6667 \ge 0{,}66$; also $k = 0$. \quad b) Der eine Wert zeigt, dass $k$ die Bedingung noch erfüllt, der andere, dass $k + 1$ sie schon verletzt. Erst beide zusammen zeigen, dass $k$ die größte Zahl ist.}
